% Consolidated 40-elementary-particles from ce5faa3e9e459c6fa8c4f0ed5dbc763b44d938dd; original TeX fragments preserved.

% === BEGIN 40-elementary-particles-291-300.tex blob 5470cc8f6890122e14e3739e4dab80f698d7a61c ===
% Irodov Part VI §6.7 original problems 291--300.
% User-supplied Arihant PDF pp.274--275 (printed pp.265--266);
% answers PDF p.360 (printed 351). Particle rest energies are MeV.
\elementaryproblem{291}{265}
质子的动量大小分别为
$0.10,\ 1.0,\ 10\,\mathrm{GeV}/c$，
求三种情形的动能。
\exsolution
质子的静止能量$m_pc^2=0.93827\,\mathrm{GeV}$。
由相对论质能关系
$E^2=p^2c^2+m_p^2c^4$、
$T=E-m_pc^2$，得到
\[
 \boxed{T(p)=\sqrt{p^2c^2+m_p^2c^4}-m_pc^2}.
\]
代入三种动量
\[
 \boxed{T_{0.10}=0.00531\,\mathrm{GeV}
 =5.31\,\mathrm{MeV}},\qquad
 \boxed{T_{1.0}=0.433\,\mathrm{GeV}},\qquad
 \boxed{T_{10}=9.11\,\mathrm{GeV}}.
\]
\excheck 本版书末首项印$0.05\,\mathrm{GeV}$，
但动量仅$0.10\,\mathrm{GeV}/c$时
非相对论检验$T\simeq p^2/(2m_p)=5.3$MeV；
书末首项有约十倍数量级差异，需记录。

\elementaryproblem{292}{265}
$\pi$介子的动能超过其静止能量$\eta=1.2$倍，
即$T=\eta m_\pi c^2$。
固有平均寿命$\tau_0=25.5\,\mathrm{ns}$，
求实验室系平均飞行距离。
\exsolution
洛伦兹因子
$\gamma=1+T/(m_\pi c^2)=1+\eta$。
平均寿命延长至$\gamma\tau_0$，
速率$v=\beta c$，
所以
\[
 \overline l=\beta\gamma c\tau_0
 =c\tau_0\sqrt{\gamma^2-1}
 =\boxed{c\tau_0\sqrt{\eta(\eta+2)}}.
\]
代入$\eta=1.2$得
\[
 \boxed{\overline l\simeq
 (299792458)(25.5\times10^{-9})\sqrt{1.2(3.2)}
 \simeq15.0\,\mathrm m}.
\]
\excheck 既要考虑寿命膨胀$\gamma$，
也要计算运动速度$\beta c$；
不可简单取$c\gamma\tau_0$。

\elementaryproblem{293}{266}
负$\pi$介子动能
$T=100\,\mathrm{MeV}$，
从产生至衰变的平均实验室距离
$l=11\,\mathrm m$。
求其固有平均寿命。
\exsolution
定义$m_\pi c^2\simeq139.57\,\mathrm{MeV}$。
由$pc=\sqrt{T(T+2m_\pi c^2)}$，
$\beta\gamma=pc/(m_\pi c^2)$。
相对论飞行距离$l=\beta\gamma c\tau_0$，
反解
\[
 \boxed{\tau_0=
 \frac{l\,m_\pi c^2}
 {c\sqrt{T(T+2m_\pi c^2)}}}.
\]
代入$T=100$MeV及$l=11$m，
\[
 \boxed{\tau_0\simeq2.63\times10^{-8}\,\mathrm s
 \simeq26\,\mathrm{ns}}.
\]
\excheck 得到的固有寿命与低速$\pi^\pm$
约26ns的寿命一致，说明已正确剥离时间膨胀。

\elementaryproblem{294}{266}
一束单能负$\pi$介子，其动能恰等于静止能量
$T=m_\pi c^2$，
在两处相距$l=20\,\mathrm m$的束流截面测通量。
固有平均寿命$\tau_0=25.5\,\mathrm{ns}$。
求后、前两截面的通量比。
\exsolution
$T=m_\pi c^2$意味着$\gamma=2$，
$\beta\gamma=\sqrt{\gamma^2-1}=\sqrt3$。
平均衰变长度
$L=\beta\gamma c\tau_0=\sqrt3c\tau_0$。
稳定束源条件下通过距离$l$的生存概率
就是通量比：
\[
 \boxed{\frac{J(l)}{J(0)}
 =\exp\!\left(-\frac{l}{\sqrt3 c\tau_0}\right)}.
\]
代入得到
$\boxed{J(l)/J(0)\simeq0.22}$。
\excheck 两个截面之间只有粒子衰变引起的损失；
不计束流几何发散或物质吸收。

\elementaryproblem{295}{266}
静止的正$\pi$介子衰变为
$\mu^+$和中微子，
求$\mu^+$动能及中微子能量。
\exsolution
取$m_\pi c^2=139.57$MeV、
$m_\mu c^2=105.66$MeV，
中微子静质量近似为零。
静止母粒子的二体衰变满足
$\boldsymbol p_\mu=-\boldsymbol p_\nu$。
由$E_\mu+E_\nu=m_\pi c^2$与
$E_\nu=p_\nu c$，
\[
 E_\nu=\frac{m_\pi^2-m_\mu^2}{2m_\pi}c^2,
 \quad
 T_\mu=\frac{(m_\pi-m_\mu)^2}{2m_\pi}c^2.
\]
代入
\[
 \boxed{T_\mu\simeq4.12\,\mathrm{MeV}},\qquad
 \boxed{E_\nu\simeq29.79\,\mathrm{MeV}}.
\]
\excheck 中微子能量与$\mu$子的动能不同，
因为$\mu$子携带显著静止能量。

\elementaryproblem{296}{266}
静止的$\Sigma^-$超子经
$\Sigma^-\to n+\pi^-$衰变，
求末态中子的动能。
\exsolution
取静止能量
$m_\Sigma c^2\simeq1197.45$MeV、
$m_nc^2\simeq939.57$MeV、
$m_\pi c^2\simeq139.57$MeV。
二体衰变的中子总能量
\[
 E_n=\frac{m_\Sigma^2+m_n^2-m_\pi^2}
 {2m_\Sigma}c^2,
\]
因此
\[
 \boxed{T_n=E_n-m_nc^2
 =\frac{(m_\Sigma-m_n)^2-m_\pi^2}
 {2m_\Sigma}c^2
 \simeq19.6\,\mathrm{MeV}}.
\]
\excheck 两粒子等动量而非等动能；
书末近似$19.5$MeV。

\elementaryproblem{297}{266}
一个初始静止的正$\mu$子衰变成
正电子和两个中微子。
求正电子可能具有的最大动能。
\exsolution
三体衰变中，当两中微子近似平行同向、
其组合的不变质量达到最小值零时，
正电子总能量最大。
视“两个中微子”为一个零质量系统，
二体运动学给
\[
 E_{e,\max}=
 \frac{m_\mu^2+m_e^2}{2m_\mu}c^2.
\]
减正电子静能
\[
 \boxed{T_{e,\max}=
 \frac{(m_\mu-m_e)^2}{2m_\mu}c^2
 \simeq52.32\,\mathrm{MeV}}.
\]
\excheck 若忽略电子质量，
近似为$m_\mu c^2/2\simeq52.8$MeV；
书末写约52.5MeV，属于所用质量常数精度差异。

\elementaryproblem{298}{266}
一个初始静止的电中性粒子
衰变为动能$T_p=5.3\,\mathrm{MeV}$的质子
及$\pi^-$。求母粒子的静止质量并指出名称。
\exsolution
静止母粒子衰变后二体动量等大反向，
因此
\[
 pc=\sqrt{T_p(T_p+2m_pc^2)},\qquad
 E_\pi=\sqrt{p^2c^2+m_\pi^2c^4}.
\]
母粒子静能
$Mc^2=m_pc^2+T_p+E_\pi$，
取$m_pc^2=938.27$MeV、
$m_{\pi^-}c^2=139.57$MeV，得
\[
 \boxed{Mc^2\simeq1115.2\,\mathrm{MeV}
 \approx1.115\,\mathrm{GeV}}.
\]
对应粒子为
$\boxed{\Lambda^0\text{ 超子}}$，
其标准静止能量约1115.7MeV。
\excheck 不能把质子动能和$\pi$的静止能直接相加；
$\pi$粒子还携带与质子等模的反冲动量。

\elementaryproblem{299}{266}
负$\pi$介子以动能$T_\pi=50\,\mathrm{MeV}$
飞行时衰变成$\mu$子和中微子。
若中微子在实验室系中与母介子运动方向
成$90^\circ$，求中微子能量。
\exsolution
四动量守恒
$P_\mu=P_\pi-P_\nu$，
结合洛伦兹不变量
$P_\mu^2=m_\mu^2c^2$。
因为中微子与母介子动量垂直，
$P_\pi\cdot P_\nu=E_\pi E_\nu/c^2$，
且$P_\nu^2=0$，所以
\[
 m_\mu^2c^4=m_\pi^2c^4-2E_\pi E_\nu,
 \qquad E_\pi=m_\pi c^2+T_\pi.
\]
解得
\[
 \boxed{E_\nu=
 \frac{(m_\pi^2-m_\mu^2)c^4}
 {2(m_\pi c^2+T_\pi)}
 \simeq21.9\,\mathrm{MeV}\approx22\,\mathrm{MeV}}.
\]
\excheck 实验室垂直出射条件使初态
空间动量点积为零，不能套静止母粒子的29.8MeV。

\elementaryproblem{300}{266}
动能$T_\Sigma=320\,\mathrm{MeV}$的
$\Sigma^+$超子飞行中衰变为
一粒电中性粒子和一粒$\pi^+$。
$\pi^+$以相对母粒子前进方向$90^\circ$出射，
其动能$T_\pi=42\,\mathrm{MeV}$。
求中性产物的静止质量。
\exsolution
由四动量
$P_X=P_\Sigma-P_\pi$，
及两动量在实验室空间中垂直，
得到
\[
 m_X^2c^4=
 m_\Sigma^2c^4+m_\pi^2c^4-2E_\Sigma E_\pi,
\]
\[
 E_\Sigma=m_\Sigma c^2+T_\Sigma,\qquad
 E_\pi=m_\pi c^2+T_\pi.
\]
取$m_{\Sigma^+}c^2\simeq1189.37$MeV、
$m_{\pi^+}c^2\simeq139.57$MeV，
\[
 \boxed{m_Xc^2\simeq
 \sqrt{1189.37^2+139.57^2-
 2(1189.37+320)(139.57+42)}
 \simeq941\,\mathrm{MeV}\approx0.94\,\mathrm{GeV}}.
\]
中性产物为$\boxed{n\text{（中子）}}$。
\excheck 输入值取两位有效数字，
所得约941MeV与中子939.6MeV在近似精度内一致。
% === END 40-elementary-particles-291-300.tex ===

% === BEGIN 40-elementary-particles-301-310.tex blob f798033b9ee66095206e8fab9af0dbe184d5c3bc ===
% Irodov Part VI §6.7, original 301--310; final 10 problems of the book.
% Source PDF pp.275--276 (printed 266--267); answers PDF p.360 (printed 351).
% Lists of reactions in problems 308--310 checked against the scanned page image.
\elementaryproblem{301}{266}
一个运动中的中性$\pi^0$介子衰变为
能量相等的两个$\gamma$光子，
两光子的传播方向夹角为
$\theta=60^\circ$。
求介子的动能以及每个光子能量。
\exsolution
记每个光子能量为$\varepsilon$。
两光子总能量$E_\pi=2\varepsilon$，
总动量模
$p_\pi=(2\varepsilon/c)\cos(\theta/2)$。
由初态介子的质能不变量
\[
 m_\pi^2c^4=E_\pi^2-p_\pi^2c^2
 =4\varepsilon^2\sin^2\frac{\theta}{2},
\]
解得
\[
 \boxed{\varepsilon=
 \frac{m_\pi c^2}{2\sin(\theta/2)}}.
\]
介子的动能
$T_\pi=2\varepsilon-m_\pi c^2$，
所以
\[
 \boxed{T_\pi=m_\pi c^2
 \left[\frac1{\sin(\theta/2)}-1\right]}.
\]
对$\theta=60^\circ$，
$\sin30^\circ=1/2$，
\[
 \boxed{T_\pi=m_\pi c^2\simeq135\,\mathrm{MeV}},\qquad
 \boxed{\varepsilon=m_\pi c^2\simeq135\,\mathrm{MeV}}.
\]
\excheck 两光子在实验室中能量相等，
但方向并非相反，所以母粒子可以携带非零动量。

\elementaryproblem{302}{266}
一个静质量为$m$、动能为$T$的粒子，
轰击初始静止的静质量为$M$的粒子。
反应生成若干静质量为
$m_1,m_2,\ldots$的新粒子。
利用不变量$E^2-p^2c^2$
证明相对论反应阈值公式。
\exsolution
实验室初态总能量
$E_{\rm in}=(mc^2+T)+Mc^2$，
三动量等于入射粒子动量$p$。
因此
\[
 s c^4=E_{\rm in}^2-p^2c^2
 =m^2c^4+M^2c^4+2M(mc^2+T)c^2,
\]
这里$s$取质量平方量纲。
阈值处末态在质心系没有相对动量，
所以末态不变质量达到
$\mathcal M_f=m_1+m_2+\cdots$，
从而$s=\mathcal M_f^2$。
解$T$得
\[
 \boxed{T_{\rm th}=
 \frac{(\sum_i m_i)^2-(m+M)^2}
 {2M}\,c^2}.
\]
\excheck 与第271题非相对论近似不同，
本式以四动量不变性处理较高能量下的末态反冲，
可用于粒子产生阈值计算。

\elementaryproblem{303}{266}
动能$T=750\,\mathrm{keV}$的正电子
撞上静止的自由电子并湮灭，
产生能量相同的两个$\gamma$光子。
求两光子夹角。
\exsolution
入射正电子总能量
$E_+=m_ec^2+T$，
初始总能量$E_{\rm tot}=2m_ec^2+T$。
两个光子各携带
$\varepsilon=E_{\rm tot}/2$。
初态不变量
\[
 s c^4=(p_++p_-)^2c^2
 =4m_e^2c^4+2m_ec^2T
 =2m_ec^2(2m_ec^2+T).
\]
末态两光子夹角为$\theta$，
则相同能量光子的不变量
$s c^4=4\varepsilon^2\sin^2(\theta/2)$。
因此
\[
 \boxed{\sin^2\frac\theta2=
 \frac{2m_ec^2}{T+2m_ec^2}},\quad
 \boxed{\theta=
 2\arcsin\sqrt{\frac{2m_ec^2}{T+2m_ec^2}}}.
\]
取$2m_ec^2=1022$keV，得
$\boxed{\theta\simeq98.8^\circ\approx99^\circ}$。
\excheck 当正电子动能$T\to0$，
两光子趋于背向发射$\theta\to180^\circ$，
验证了角度公式的极限行为。

\elementaryproblem{304}{266}
求入射$\gamma$光子与静止靶粒子反应时的阈值：
(a) 在自由电子场中产生一对$e^-e^+$；
(b) 在静止质子场中产生一对带相反电荷的$\pi$介子。
\exsolution
质量为$M$的静止靶吸收高能光子后，
初态不变量为$M^2c^4+2Mc^2\varepsilon$。
在阈值处末态总静质量为$\mathcal M_f$，
故
$\varepsilon_{\rm th}=
[(\mathcal M_f)^2-M^2]c^2/(2M)$。
(a) 反应
$\gamma+e^-\to e^-+e^-+e^+$，
末态总质量$3m_e$，故
\[
 \boxed{\varepsilon_{\rm th}^{(a)}
 =\frac{(3m_e)^2-m_e^2}{2m_e}c^2
 =4m_ec^2\simeq2.04\,\mathrm{MeV}}.
\]
(b) 反应
$\gamma+p\to p+\pi^++\pi^-$，
末态总质量$m_p+2m_\pi$，故
\[
 \boxed{\varepsilon_{\rm th}^{(b)}=
 \frac{(m_p+2m_\pi)^2-m_p^2}{2m_p}c^2
 =2m_\pi c^2\!\left(1+\frac{m_\pi}{m_p}\right)
 \simeq321\,\mathrm{MeV}}.
\]
\excheck 对自由电子靶的对产生阈值并非单纯
$2m_ec^2$，因为初态单光子有动量，
末态必须分走反冲动能。

\elementaryproblem{305}{266}
质子束轰击静止氢靶，求以下反应的
最小质子动能：
(a) $p+p\to p+p+p+\bar p$；
(b) $p+p\to p+p+\pi^0$。
\exsolution
对静止质子靶使用第302题阈值公式，
入射、靶质量都为$m_p$。
(a) 末态四粒子的总静质量为$4m_p$，
所以
\[
 \boxed{T_{{\rm th},a}=
 \frac{(4m_p)^2-(2m_p)^2}{2m_p}c^2
 =6m_pc^2\simeq5.63\,\mathrm{GeV}}.
\]
(b) 末态总静质量$2m_p+m_{\pi^0}$，
所以
\[
 \boxed{T_{{\rm th},b}=
 \frac{(2m_p+m_{\pi^0})^2-(2m_p)^2}{2m_p}c^2
 =2m_{\pi^0}c^2+\frac{m_{\pi^0}^2c^2}{2m_p}
 \simeq0.280\,\mathrm{GeV}}.
\]
\excheck (a) 产生反质子时同时必须有一个额外质子，
以保证总重子数从初态2到末态2。
原书书末第305(a)题的结果末尾印为“$5.6\,\mathrm{MeV}$”，
但由严格阈值$6m_pc^2$可知必为约
$5.6\,\mathrm{GeV}$，相差$10^3$；
这里明确保留并指出原书单位的错误。

\elementaryproblem{306}{267}
$\pi$介子束轰击静止氢靶，
求产生下列粒子的入射动能阈值：
(a) $\pi^-+p\to K^++\Sigma^-$；
(b) $\pi^0+p\to K^++\Lambda^0$。
\exsolution
对入射质量$m_\pi$、静止质子靶$m_p$，
阈值条件
\[
 T_{\rm th}=
 \frac{(M_f)^2-(m_\pi+m_p)^2}{2m_p}c^2.
\]
(a) 取
$m_{\pi^\pm}c^2=139.57$MeV、
$m_pc^2=938.27$MeV、
$m_{K^+}c^2=493.68$MeV、
$m_{\Sigma^-}c^2=1197.45$MeV，
\[
 \boxed{T_{{\rm th},a}\simeq0.905\,\mathrm{GeV}
 \approx0.90\,\mathrm{GeV}}.
\]
(b) 取
$m_{\pi^0}c^2=134.98$MeV、
$m_{\Lambda^0}c^2=1115.68$MeV，
\[
 \boxed{T_{{\rm th},b}\simeq0.766\,\mathrm{GeV}
 \approx0.77\,\mathrm{GeV}}.
\]
\excheck 末态两个粒子在阈值处
具有相同的质心速度，不是实验室系静止；
应使用质量平方不变量公式。

\elementaryproblem{307}{267}
某电中性基本粒子具有
同位旋第三分量$T_z=+1/2$，
重子数$B=+1$。
求其奇异数$S$、超荷$Y$，
并指出粒子名称。
\exsolution
强子满足盖尔曼--西岛关系
\[
 Q=T_z+\frac{Y}{2},\qquad Y=B+S.
\]
由$Q=0,T_z=+1/2$得$Y=-1$，
再由$B=1$得$S=Y-B=-2$。
因此
\[
 \boxed{S=-2,\qquad Y=-1},\qquad
 \boxed{\text{粒子为中性级联超子 }\Xi^0}.
\]
\excheck $\Xi^0$含两个奇异夸克，
故总奇异数恰为$-2$，
与同位旋投影和电荷关系一致。

\elementaryproblem{308}{267}
判定下列过程哪些违反轻子数守恒
（记$\nu_e,\nu_\mu$为各自中微子）：
\[
\begin{array}{ll}
(1)\ n\to p+e^-+\nu_e;&
(2)\ \pi^+\to\mu^++e^-+e^+;\\
(3)\ \pi^-\to\mu^-+\nu_\mu;&
(4)\ p+e^-\to n+\nu_e;\\
(5)\ \mu^+\to e^++\nu_e+\bar\nu_\mu;&
(6)\ K^-\to\mu^-+\bar\nu_\mu.
\end{array}
\]
\exsolution
约定每种轻子的轻子数为$+1$、
相应反粒子为$-1$；
介子、核子和光子轻子数为零。
逐条对电子/缪子轻子数分别计数：
(1) 末态$e^-$及$\nu_e$都为$+1$，
初态为0，禁戒；允许的$\beta^-$衰变应发
$\bar\nu_e$而不是$\nu_e$。
(2) $e^-+e^+$轻子数互消，
但单独出现$\mu^+$的$L_\mu=-1$，
禁戒。
(3) $\mu^-$和$\nu_\mu$均有$L_\mu=+1$，
初态0，禁戒；应改为$\bar\nu_\mu$。
(4) 初态电子轻子数$+1$，
末态$\nu_e$为$+1$，允许。
(5) 左端$\mu^+$有$L_\mu=-1$，
右端$\bar\nu_\mu$亦为$-1$；
$e^+$及$\nu_e$相消，允许。
(6) $\mu^-$与$\bar\nu_\mu$
分别为$+1,-1$，允许。
\[
 \boxed{\text{违反轻子数守恒的过程为(1)、(2)、(3)。}}
\]
\excheck 第(1)与第(3)题面\emph{明确写的是
中微子而非反中微子}，不能把符号看漏；
本版书末与判定一致。

\elementaryproblem{309}{267}
判断下列过程哪些违反强相互作用的奇异数守恒：
\[
\begin{array}{ll}
(1)\ \pi^-+p\to\Sigma^-+K^+;&
(2)\ \pi^-+p\to\Sigma^++K^-;\\
(3)\ \pi^-+p\to K^++K^-+n;&
(4)\ n+p\to\Lambda^0+\Sigma^+;\\
(5)\ \pi^-+n\to\Xi^-+K^++K^-;&
(6)\ K^-+p\to\Omega^-+K^++K^0.
\end{array}
\]
\exsolution
奇异数分配：
$\pi,n,p$为0，
$K^+,K^0$为$+1$，
$K^-$为$-1$，
$\Sigma,\Lambda$为$-1$，
$\Xi$为$-2$，$\Omega$为$-3$。
逐项比较初、末态总奇异数：
\[
\begin{array}{c|c|c}
\text{过程}&S_{\rm in}&S_{\rm out}\\\hline
(1)&0&-1+1=0\\
(2)&0&-1-1=-2\\
(3)&0&1-1+0=0\\
(4)&0&-1-1=-2\\
(5)&0&-2+1-1=-2\\
(6)&-1&-3+1+1=-1
\end{array}
\]
所以
\[
 \boxed{\text{违反强相互作用奇异数守恒的为
 (2)、(4)、(5)。}}
\]
\excheck 过程(6)初始已经携带$S=-1$，
只要求末态总奇异数同为$-1$，
而非必须等于零。

\elementaryproblem{310}{267}
指出以下六个过程为何被禁戒：
\[
\begin{array}{ll}
(1)\ \Sigma^-\to\Lambda^0+\pi^-;&
(2)\ \pi^-+p\to K^++K^-;\\
(3)\ K^-+\pi^-\to\Omega^-+K^++K^0;&
(4)\ n+p\to\Sigma^++\Lambda^0;\\
(5)\ \pi^-\to\mu^-+e^++e^-;&
(6)\ \mu^-\to e^-+\nu_e+\bar\nu_\mu.
\end{array}
\]
\exsolution
(1) 能量守恒/静质量阈值：
$m_{\Sigma^-}c^2\simeq1197$MeV，
$m_{\Lambda^0}c^2+m_{\pi^-}c^2
\simeq1116+140=1256$MeV；
静止母粒子无法产生静质量更大的产物。
(2) 重子数守恒：
初态$p$有$B=+1$，末态仅为介子$B=0$。
(3) \emph{电荷和重子数同时不守恒}：
初态$K^-+\pi^-$的电荷为$-1-1=-2$，
重子数$B_{\rm in}=0$；
末态$\Omega^-+K^++K^0$的电荷为$-1+1+0=0$，
但还出现了一个$\Omega^-$重子，
使$B_{\rm out}=1$。
因此即使人为改变运动能量，
本过程仍违反两条独立的严格守恒律。
(4) 奇异数守恒：
初态$n+p$的$S=0$，
末态$\Sigma^++\Lambda^0$的$S=-2$，
不能经守恒奇异数的强相互作用产生。
(5) 轻子数守恒：
初态$\pi^-$的轻子数为0，
末态$\mu^-$有缪轻子数$+1$，
而$e^+e^-$电子轻子数互消。
(6) 分味轻子数守恒：
左端$\mu^-$有$L_\mu=+1$，
右端$\bar\nu_\mu$有$L_\mu=-1$，
而$e^-+\nu_e$使$L_e=+2$。
虽然总轻子数在某些计数下形式相同，
电子轻子数与缪轻子数分别不守恒。
因此
\[
 \boxed{\begin{array}{l}
 (1)\ \text{能量不足};\quad(2)\ \text{重子数};\quad
 (3)\ \text{电荷与重子数};\\
 (4)\ \text{奇异数};\quad(5)\ \text{轻子数};\quad
 (6)\ \text{电子与缪轻子数各自不守恒}.
 \end{array}}
\]
\excheck 原书最后一题的六项以扫描图片逐项核对；
第(3)项同时违反电荷和重子数守恒，
只报告其中一条会遗漏独立的禁戒原因。
第(6)项特别需要分开统计电子与缪子轻子数，
不可只看两边总轻子数是否巧合相等。
% === END 40-elementary-particles-301-310.tex ===
